\originalchapter{白矮星冷却、简并压力与致密星极限}
本章对应8.901第14--15和18次课。第\ref{ch:compact}章已给零温费米气体和Lane--Emden推导；这里连接质量半径律、冷却时钟及GR平衡。Deepto Chakrabarty的课程课历将中子星／脉冲星课列为David Kaplan客座；没有公开该堂完整笔记，以下为独立教学补写。

\originalsection{从电子数到质量极限}
零温两自旋电子填满费米球，$p_F=\hbar(3\pi^2n_e)^{1/3}$。压强是各向同性动量通量，$P=(3\pi^2\hbar^3)^{-1}\int_0^{p_F}p^3v(p)\dd p$。用$\rho=\mu_em_un_e$，
\[
 K_{\rm NR}=\frac{\hbar^2(3\pi^2)^{2/3}}{5m_e(\mu_em_u)^{5/3}},
 \quad P_{\rm NR}=K_{\rm NR}\rho^{5/3},
\]
\[
 K_{\rm UR}=\frac{\hbar c(3\pi^2)^{1/3}}{4(\mu_em_u)^{4/3}},
 \quad P_{\rm UR}=K_{\rm UR}\rho^{4/3}.
\]
这里$m_u$为核子质量归一；粗略用$m_p$的结构估计会有小系数差，核质量表不能照此替换。完全电离碳氧$\mu_e\simeq2$，但$\mu$还数离子，两者不能混用。

多方$n=3/2$质量半径缩放给$R\propto K_{\rm NR}M^{-1/3}/G$：更重白矮星更小。相对论电子的$n=3$族则
\[
 M=4\pi\left(\frac{K_{\rm UR}}{\pi G}\right)^{3/2}w_3,\qquad
 w_3=-\xi_1^2\theta'(\xi_1).
\]
中心密度完全取消，形成质量尺度；$w_3$来自实际Lane--Emden积分，PS4.3给2.01824。代常数约$5.83\mu_e^{-2}\Ms$，$\mu_e=2$约$1.46\Ms$。常用较精细的约1.4太阳质量还涉及非理想和演化修正，不能把本理想模型小数当作所有白矮星的精确失稳质量。

同质绝热缩放模型的稳定要求$\Gamma>4/3$（第\ref{ch:stellar}章），电子由$5/3$渐向$4/3$提供了极限的物理方向。真实全星EOS连续过渡，不是将整颗白矮星突然切换两个指数；电子俘获、组成、库仑作用、自转及GR都会影响终点。

\originalsection{白矮星发光靠冷却，压强不靠同一温度}
简并电子支持星体而离子可储热，散失热量不一定像理想气体主序核心那样使供压崩溃。高导热核心近等温，薄非简并包层将核心温度$T_c$连接到光度。若离子比热近常数，$C\dd T_c/\dd t=-L(T_c)$；在Mestel型$L\propto MT_c^{7/2}$近似下，积分给晚时$t\propto T_c^{-5/2}\propto L^{-5/7}$。这依赖包层不透明度、组成和EOS匹配，不能只知质量和表面温度就假定唯一时钟。

离子Coulomb耦合参数$\Gamma=Z^2e^2/(a_{\rm ion}k_BT_c)$（Gaussian单位），$a_{\rm ion}=(3/(4\pi n_{\rm ion}))^{1/3}$。温度降低使$\Gamma$增大，结晶可释放潜热、改变热容和组成分离，延迟冷却。PS6.2使用题给$\Gamma=160$及一个明确附加包层模型；它算出的条件冷却尺度没有被宣称为真实白矮星年龄。

白矮星光度函数的低光度端可约束最早形成星的冷却时间，再加前身主序寿命，估星族最低年龄。需要选择、质量分布和消光校正；“最冷一颗白矮星”的温度不是直接等于银河系年龄。PS6.3沿原题时钟参数并说明其平均粒子质量约定的物理差别。

\originalsection{中子星需把压强本身的引力纳入}
更高密度下电子俘获、核相互作用及中子简并改变EOS；不能把电子白矮星模型仅换成中子质量就视为真实中子星。球对称静态TOV方程
\[
 \frac{\dd P}{\dd r}=
 -\frac{G[m+4\pi r^3P/c^2](\rho+P/c^2)}
 {r^2[1-2Gm/(rc^2)]},\quad
 \frac{\dd m}{\dd r}=4\pi r^2\rho
\]
中的$\rho=\epsilon/c^2$是总能量密度。三个GR修正分别涉及压强的质量贡献、惯性密度和曲率分母，均会使同一密度的压强支持更困难；弱场、$P\ll\rho c^2$时恢复Newton式。

EOS必须给$P(\epsilon)$，因果性要求$0\le\dd P/\dd\epsilon\le1$。PS8.2照原题粗略重子密度换算实际求出两族TOV曲线，也发现给定线性插值超声速；这种模型最大质量不能据为观测上限。中心正规展开、表面条件、收紧容差和能量单位的核查都是解释数值图所需的内容。

均匀密度解析球的有限中心压强要求$R>(9/8)(2GM/c^2)$，PS8.1完整证明；其不可压EOS本身不因果。不同概念必须分开：事件视界半径、静态压力极限、某EOS的质量转折、自转质量极限和稳定性判据不是同一条边界。
